\documentclass[11pt,a4paper]{article}

% =============================================================================
% MINT + PM Web Scraper -- Technical Documentation
% Paste this entire file into a new Overleaf project as main.tex and
% compile with pdfLaTeX.
% =============================================================================

\usepackage[margin=2.5cm]{geometry}
\usepackage[T1]{fontenc}
\usepackage[utf8]{inputenc}
\usepackage{lmodern}
\usepackage{booktabs}
\usepackage{longtable}
\usepackage{array}
\usepackage{xcolor}
\usepackage{listings}
\usepackage{hyperref}
\usepackage{parskip}
\usepackage{enumitem}

\setlist{nosep, leftmargin=1.4em}

\hypersetup{
    colorlinks=true,
    linkcolor=blue!60!black,
    urlcolor=blue!60!black,
    pdftitle={MINT + PM Web Scraper - Technical Documentation},
    pdfauthor={Bobrick Washroom Equipment - Maintenance Engineering}
}

\definecolor{codebg}{RGB}{245,245,245}
\lstset{
    basicstyle=\ttfamily\small,
    backgroundcolor=\color{codebg},
    frame=single,
    framerule=0pt,
    breaklines=true,
    columns=fullflexible,
    keepspaces=true,
    showstringspaces=false
}

\newcommand{\code}[1]{\texttt{#1}}

\title{\textbf{MINT + PM Web Scraper}\\[4pt]
       \Large Technical Documentation\\[4pt]
       \normalsize Bobrick Maintenance Dashboard \& Preventive-Maintenance Data Platform}
\author{Bobrick Washroom Equipment --- Maintenance Engineering}
\date{\today}

\begin{document}

\maketitle
\tableofcontents
\newpage

% =============================================================================
\section{Introduction}
% =============================================================================

\subsection{Purpose}
This document describes the architecture, data model, APIs, and operational
procedures of \textbf{MINT} (Maintenance INTelligence) and its companion
\textbf{PM Web Scraper}. Together they form a self-hosted maintenance
management platform for Bobrick Washroom Equipment that:

\begin{itemize}
    \item Scrapes historical work-order data from the legacy Bobrick PM
          system (\code{circaweb.bobrick.com/PME}) using a live, logged-in
          Chrome browser session, across multiple company divisions.
    \item Serves a single-page web dashboard with a four-level drill-down:
          \emph{Company} $\rightarrow$ \emph{Division} $\rightarrow$
          \emph{Department} $\rightarrow$ \emph{Machine}.
    \item Acts as the system of record going forward: users create and edit
          work orders, machines, departments, vendor contacts, spare parts
          and machine information directly in MINT, layered on top of the
          frozen scraped snapshot.
    \item Generates and maintains AI-assisted per-machine troubleshooting
          checklists (LLM-backed, with full version history and manual-edit
          preservation).
    \item Sends e-mail notifications for new unscheduled work orders,
          completed work orders, and a weekly scheduled-maintenance digest.
\end{itemize}

\subsection{Scope}
The system runs as a single Python/Flask process, reachable on the local
LAN. All state lives on the local file system: JSON snapshots, SQLite
databases, Markdown checklists, and uploaded attachments. The only external
services are optional LLM providers (Ollama Cloud or Google Gemini) and an
optional SMTP server for notifications.

\subsection{Terminology}
\begin{description}
    \item[PM system] The legacy Bobrick preventive-maintenance web
          application being scraped.
    \item[MINT] The dashboard application described in this document.
    \item[WO] Work order. \emph{Unscheduled} (breakdown/repair) or
          \emph{Scheduled} (preventive maintenance).
    \item[Division] A Bobrick business unit (e.g.\ BLA, BED, BMC). Each is
          scraped and stored independently.
    \item[EQ ID] Equipment identifier; machines are addressed by the numeric
          portion of their equipment ID.
    \item[Machine / Equipment] Used interchangeably throughout the codebase.
    \item[Checklist / Guide] The per-machine AI troubleshooting checklist
          (Markdown, stored under \code{guides/}).
\end{description}

% =============================================================================
\section{System Overview}
% =============================================================================

\subsection{High-Level Architecture}
The platform is organized into five logical layers.

\begin{enumerate}
    \item \textbf{Acquisition layer} --- \code{scraper.py},
          \code{chrome\_session.py}, \code{run\_parallel.py},
          \code{nightly\_update.py}. Drives a real, already-logged-in Chrome
          browser via the DevTools remote-debugging protocol to extract
          every work order, comment, and attachment from the PM system.
          Output is written to per-department JSON/CSV snapshot files, and
          is division-aware so any Bobrick division can be scraped with the
          same pipeline.
    \item \textbf{Storage layer} --- Frozen scraped JSON snapshots
          (read-only) plus two SQLite databases: \code{mint\_data/mint.db}
          for all user-created and user-edited data, and
          \code{guides/versions.db} for checklist version history. Uploaded
          files live under \code{mint\_data/attachments/}.
    \item \textbf{Application layer} --- \code{server.py}, a Flask
          application that merges the frozen snapshot with the live store at
          load time, exposes a JSON REST API, and hosts background jobs
          (weekly e-mail digest, recurring-PM generation, Chrome session
          bootstrap).
    \item \textbf{Intelligence layer} --- \code{analyze\_equipment.py} and
          \code{guide\_engine.py}. LLM-backed trend analysis, monthly
          synopses, and troubleshooting-checklist generation via Ollama
          Cloud or Google Gemini.
    \item \textbf{Presentation layer} --- \code{webapp/index.html}, a single
          self-contained single-page application built with React (via
          \code{htm} and Babel standalone), Tailwind CSS, \code{marked} for
          Markdown rendering, and \code{jsPDF} for client-side PDF export.
          All frontend dependencies are vendored locally so the app works
          fully offline on the LAN.
\end{enumerate}

\subsection{Data Flow}
\begin{lstlisting}
 Legacy PM site (circaweb.bobrick.com/PME), any division
        |  logged-in Chrome, DevTools port auto-detected
        v
 scraper.py / run_parallel.py / nightly_update.py
        |  writes
        v
 divisions/<slug>/work_orders_*.json/.csv   (frozen snapshot, BLA at root)
        |                                mint_data/mint.db  (user writes)
        |                                guides/*.md + versions.db
        v                                /
 server.py  reload_data(): merge snapshot + store, apply overrides,
        |   hide inactive items, index by wo_id
        v
 JSON REST API (/api/...)  <---->  webapp/index.html (React SPA)
        |
        v
 mint_email.py  (SMTP notifications, inert until configured)
\end{lstlisting}

\subsection{Key Design Principles}
\begin{description}
    \item[Frozen snapshot + write layer] Scraped PM data is never mutated.
          User changes live in SQLite and are merged over the snapshot at
          load time. Field-level edits to any WO (scraped or manual) are
          stored as JSON patches.
    \item[Attributability] Every mutating API call requires an
          \code{author} (``Your name'') and writes an audit-log row.
    \item[Soft deletion] Departments and machines are never hard-deleted;
          they are flagged inactive and can be restored.
    \item[Manual-edit preservation] Operator edits to AI checklists are
          diffed into an append-only edit log that is re-injected into every
          future LLM prompt, so edits survive full regeneration.
    \item[Fail-safe scraping] Nightly rescrapes write to temp files and are
          only promoted over live data when the result looks valid; an
          expired login can never wipe good data.
    \item[Offline-first frontend] All JavaScript/CSS dependencies are
          vendored; the app requires no CDN or internet connection to
          render.
\end{description}

% =============================================================================
\section{Repository Layout}
% =============================================================================

\begin{longtable}{@{}>{\ttfamily}p{0.40\textwidth}p{0.55\textwidth}@{}}
\toprule
\textbf{Path} & \textbf{Description} \\
\midrule
\endhead
server.py & Flask backend, REST API, background schedulers. \\
mint\_store.py & SQLite write store for all user-created data. \\
mint\_email.py & SMTP notification module (inert until configured). \\
version\_store.py & Append-only checklist version history (SQLite). \\
guide\_engine.py & Checklist create/edit/archive/regenerate engine. \\
analyze\_equipment.py & LLM trend analysis, stats, prompt building, transport. \\
scraper.py & Live-browser PM work-order scraper. \\
chrome\_session.py & Persistent logged-in debug-Chrome session manager. \\
run\_parallel.py & Multi-department scrape orchestration, per division. \\
nightly\_update.py & Full rescrape + checklist regeneration cycle. \\
parse\_html.py & HTML parsing helpers for captured pages / equipment lists. \\
enrich\_assigned\_to.py & One-time backfill of Owner / Assigned-To fields. \\
build\_vendor\_contacts.py & Builds the vendor contact list from source CSVs. \\
import\_machine\_longevity.py & Imports the machine-longevity CSV into machine\_info. \\
webapp/index.html & The entire SPA frontend (single file). \\
webapp/vendor/ & Vendored react, react-dom, htm, babel, tailwind, marked, jspdf. \\
guides/ & Live checklists, baselines/, edits/, archive/, versions.db. \\
mint\_data/ & mint.db, attachments/, email\_recipients.json (git-ignored). \\
divisions/<slug>/ & Per-division scraped snapshots, pages mirror, profiles. \\
pages/ & Offline mirror of every scraped BLA page. \\
work\_orders\_*.json/.csv & Frozen BLA scraped snapshots (master + per department). \\
equipment\_data.json/.csv & Scraped equipment master list (BLA). \\
start\_chrome\_debug.bat & Manual launcher for the debug Chrome profile. \\
requirements.txt & Python dependencies. \\
.env & Local secrets/config (git-ignored; see \S\ref{sec:config}). \\
\bottomrule
\end{longtable}

% =============================================================================
\section{Acquisition Layer (Scraping)}
% =============================================================================

\subsection{Chrome Session Management}
The scraper attaches to a real Chrome browser that a human has logged into
the PM site, so no credentials are stored in code.

\begin{enumerate}
    \item Chrome is launched once with a dedicated profile and
          \code{--remote-debugging-port=0}, letting Chrome pick a free port.
    \item The chosen port is read from the \code{DevToolsActivePort} file
          Chrome writes into the profile directory, then verified against
          the debug endpoint.
    \item The window is left open so a human can re-login if the session
          expires; the same session is reused for every scrape.
\end{enumerate}

Standalone usage: \code{python chrome\_session.py} (launch and print port)
or \code{python chrome\_session.py --status} (report port and login state).

\subsection{Work-Order Scraper}
For every equipment record, the scraper:
\begin{enumerate}
    \item Opens the Equipment Dashboard for that machine.
    \item Iterates every row of the Unscheduled Work Orders grid, opening
          the Edit Work Order dialog for each.
    \item Captures date notified, urgency, problem, status, material cost,
          labor time, work performed by, downtime hours, completed
          date/time, the full comment log with timestamps, and all
          attachment filenames. Two distinct people/teams are captured on
          every WO: \code{owner} (who reported it) and \code{assigned\_to}
          (the skill/team it is assigned to) --- these are different
          concepts and are kept separate end to end.
    \item Repeats the same for the Scheduled (PM) grid.
\end{enumerate}

Every visited page is mirrored into \code{pages/} (form state baked in,
scripts stripped), enabling a full offline replay with
\code{python scraper.py --from-html} against a private headless Chrome ---
useful for testing scraper changes without touching the live site.
\code{enrich\_assigned\_to.py} is a one-time backfill script that
re-derives the owner / assigned-to fields from saved HTML exports for
records scraped before the distinction existed.

\subsection{Multi-Division Scraping}
\label{sec:divisions}
The scraper and orchestrator are division-aware so the same pipeline used
for BLA can target any division exposed by the PM system's division
dropdown (BLA, BED, BMC, BWEC, KKP, BGD, CIT):

\begin{itemize}
    \item \code{scraper.py --division <slug>} routes all output paths under
          \code{divisions/<slug>/} instead of the repository root. BLA,
          with no \code{--division} flag, is unchanged and still writes to
          the root.
    \item The Chrome session actively sets the division selector before
          scraping, since it does not inherit whatever division was last
          selected interactively.
    \item \code{--equipment-list-only} builds just that division's
          equipment master; a heuristic warns if scraped equipment IDs
          collide heavily with BLA's, indicating the selector did not take
          effect.
    \item \code{run\_parallel.py --division <slug>} discovers that
          division's departments and scrapes them the same way as BLA
          (per-department parallel jobs), namespaced under
          \code{divisions/<slug>/}.
    \item \textbf{Status:} BLA is fully scraped at the repository root. BED
          has a complete scrape under \code{divisions/bed/}. Other
          divisions are scraped incrementally with the same pipeline;
          department keys are namespaced per division (e.g.\
          \code{bed\_maintenance}) to avoid collisions with identically
          named departments in other divisions.
\end{itemize}

\subsection{Nightly Update Cycle}
One cycle (\code{python nightly\_update.py}) performs, in order:
\begin{enumerate}
    \item Archives every current checklist into a timestamped folder.
    \item Snapshots current unscheduled WO IDs per machine, to detect what
          is new after the scrape.
    \item Rescrapes all departments sequentially through the single logged
          in Chrome. Results go to temp files first and are only promoted
          over live per-department files when valid. A per-department
          timeout kills hung departments so the run continues.
    \item Merges per-department files into the master snapshot files.
    \item Regenerates checklists only for machines that gained new
          unscheduled WOs, preserving all manual edits.
    \item Reloads the server's in-memory data so the frontend updates
          immediately.
\end{enumerate}

% =============================================================================
\section{Storage Layer}
% =============================================================================

\subsection{Frozen Scraped Snapshots}
Scraped data lives in flat JSON/CSV files per department, e.g.\
\code{work\_orders\_unscheduled\_toilet\_partitions.json} (or under
\code{divisions/<slug>/} for non-BLA divisions). These files are treated as
read-only historical truth; the application never rewrites them outside the
nightly promotion step.

\subsection{MINT Write Store (SQLite)}
\label{sec:store}
Database: \code{mint\_data/mint.db}. A module-level write lock serializes
writes so ID assignment stays atomic. Manual work orders receive the prefix
\code{M-} so they can never collide with the purely numeric PM WO numbers.

\begin{longtable}{@{}>{\ttfamily}p{0.24\textwidth}p{0.71\textwidth}@{}}
\toprule
\textbf{Table} & \textbf{Purpose / key columns} \\
\midrule
\endhead
divisions & Company $\rightarrow$ division layer: \code{key} (PK), \code{name}. \\
departments & User-added departments: \code{key} (PK), \code{name}, \code{label}, \code{division\_key}. Scraped departments are configured statically. \\
machines & User-added machines: \code{eq\_id}, \code{equipment\_name}, \code{dept\_key}, \code{make}, \code{model}, \code{vendor}, \code{asset\_num}. \\
work\_orders & Manually created WOs: \code{wo\_id} (PK, \code{M-} prefixed), \code{wo\_type}, \code{department\_key}, equipment fields, full record as JSON. \\
wo\_overrides & Field-level edits applied to any WO (scraped or manual): \code{wo\_id} (PK) $\rightarrow$ JSON patch. \\
solutions & Append-only Solution log entries per WO (used instead of comments on unscheduled WOs). \\
attachments & Uploaded file metadata; bytes live under \code{mint\_data/attachments/}. \\
inactive\_items & Soft-delete flags. PK is (\code{item\_type}, \code{item\_key}). \\
audit\_log & Who did what and when: \code{wo\_id}, \code{author}, \code{action}, \code{detail}, timestamp. \\
calendar\_events & User calendar entries (date, department, equipment, title, description). \\
chart\_events & Global timeline markers rendered on every trend chart. \\
machine\_info & Editable machine information / longevity data, PK (\code{dept\_key}, \code{eq\_id}): category, location, serial/asset number, year, condition, service status, replacement cost/year, comments, LLM-parsed summary. \\
technicians & Individually tracked technician roster with aliases and an active flag. \\
roles & Shared ``assigned-to'' identities (e.g.\ ``QA Technician'') multiple people can sign in as for My Schedule, without individual tracking: \code{name} (PK), aliases, active flag, assigned \code{view\_key}. \\
view\_profiles & Named, saved sets of enabled UI features that can be assigned to a role, so it sees only a curated subset of MINT. \\
vendor\_contacts & Editable Vendor \& Utilities list: company, contact, address, phone/cell/fax, email, type, service-contract fields, linked equipment IDs. \\
vendor\_contact\_types & Contact-type lookup; deletion blocked while a contact still references the type. \\
floorplan\_items & Interactive floor-plan boxes per department: id, dept\_key, optional machine link, label, position/size, z-order. \\
machine\_nicknames & Admin-set display name overriding the scraped PM equipment name, PK (\code{dept\_key}, \code{eq\_id}). \\
spare\_parts & Division-scoped spare-parts inventory: description, digital ID, picture, primary machine link, part type, location, quantity, condition, brand, set membership, purchasing info, barcode. \\
spare\_part\_machines & Many-to-many link from a part to additional machines beyond its primary one. \\
\bottomrule
\end{longtable}

Schema creation is idempotent (\code{CREATE TABLE IF NOT EXISTS}) with
lightweight migrations for columns added after initial release.

\subsection{Checklist Version Store}
Database: \code{guides/versions.db}. Every change to a checklist --- AI
generated, AI updated, or hand-edited --- is recorded as a new immutable
version containing the full Markdown, with source (\code{ai\_generate},
\code{ai\_update}, \code{manual\_edit}, \code{import}, \code{restore}),
author, timestamp, and linked previous/next version references.

% =============================================================================
\section{Application Layer (server.py)}
% =============================================================================

\subsection{Startup Sequence}
Running \code{python server.py}: loads \code{.env}, performs the initial
data merge, best-effort launches the debug Chrome
(\code{CHROME\_AUTOSTART=1} by default), starts the weekly digest scheduler
thread, and runs Flask threaded on \code{SERVER\_HOST:SERVER\_PORT}
(default \code{0.0.0.0:5000}).

\subsection{Data Merge (reload\_data)}
The core of the read path. It atomically rebuilds the in-memory caches:
\begin{enumerate}
    \item Refreshes the technician/role lists and rebuilds the live
          department map (static scraped departments plus user-added ones,
          minus inactive ones).
    \item Auto-generates due occurrences for recurring scheduled-WO series.
    \item Loads each department's scheduled/unscheduled JSON, tagging
          records with WO type and department key.
    \item Injects manually created WOs from the store, normalizing dates so
          they sort correctly.
    \item Builds a global WO index, applies field-level overrides, and
          tags solution/attachment counts.
    \item Filters out soft-deleted machines and their WOs, then swaps all
          caches in atomically.
\end{enumerate}

\subsection{Recurring Scheduled Work Orders}
Scheduled WOs with a frequency form a series. Occurrence generation is
purely date-driven: on every reload, occurrences are created on a fixed
cadence from the seed's first due date through today plus the next
upcoming one, independent of completion status. A series can be
permanently stopped (admin-gated) while keeping existing occurrences.

\subsection{Access Control}
Independent shared passwords, all supplied via header or JSON body and all
optional (a blank password leaves that gate open so a fresh install is
never locked out):

\begin{longtable}{@{}p{0.16\textwidth}>{\ttfamily}p{0.25\textwidth}p{0.51\textwidth}@{}}
\toprule
\textbf{Gate} & \textbf{Env / header} & \textbf{Protects} \\
\midrule
\endhead
Edit & EDIT\_PASSWORD / X-Edit-Password & Checklist generate/edit/update, vendor-contact and chart-event mutations. \\
Admin & ADMIN\_PASSWORD / X-Admin-Password & Deleting work orders, viewing WO change history, stopping recurrence, spare-parts mutations, nicknames, floor-plan edits. \\
Sean (supervisor) & SEAN\_PASSWORD / X-Sean-Password & Work-order assignment and the completed-WO team dashboard. Defaults to \code{sean}. \\
Master & MASTER\_PASSWORD / X-Master-Password & Managing view profiles and assigning them to roles. Defaults to \code{mintmaster}. \\
\bottomrule
\end{longtable}

Every mutation also requires an \code{author} name, persisted to the audit
log (the frontend remembers it in \code{localStorage}).

\subsection{API Reference}
All endpoints return JSON. Read endpoints are unauthenticated on the LAN.

\subsubsection{Navigation and Dashboards}
\begin{longtable}{@{}p{0.16\textwidth}>{\ttfamily\small}p{0.45\textwidth}p{0.31\textwidth}@{}}
\toprule
\textbf{Method} & \textbf{Path} & \textbf{Description} \\
\midrule
\endhead
GET & / & Serves the SPA. \\
GET & /api/company & Company KPIs + one card per division. \\
GET & /api/division, /api/divisions/<key> & Division dashboard: departments, KPIs, inactive departments. \\
GET & /api/departments/<dept\_key> & Department dashboard: machines, KPIs, inactive machines. \\
GET & /api/departments/<dk>/machines/<eq\_id> & Machine detail. \\
GET & /api/workorders & All departments' WOs, compact, newest first. \\
GET & /api/workorder/<wo\_id> & Full work-order detail. \\
GET & /api/weekly, /api/monthly & Weekly and monthly schedule views. \\
GET & /api/completed-counts, /api/team-stats & Completed-WO counts and per-technician stats. \\
GET & /api/inactive & All soft-deleted departments/machines. \\
\bottomrule
\end{longtable}

\subsubsection{Work-Order Lifecycle}
\begin{longtable}{@{}p{0.17\textwidth}>{\ttfamily\small}p{0.44\textwidth}p{0.31\textwidth}@{}}
\toprule
\textbf{Method} & \textbf{Path} & \textbf{Description} \\
\midrule
\endhead
POST & /api/workorders & Create a WO; fires the new-unscheduled e-mail. \\
PATCH / PUT & /api/workorder/<wo\_id> & Field-level edit of any WO; fires the completed e-mail on status transition. \\
DELETE & /api/workorder/<wo\_id> & Delete a manual WO (admin-gated; scraped WOs cannot be deleted). \\
POST & /api/workorder/<wo\_id>/stop-recurrence & Stop a recurring series (admin-gated). \\
POST & /api/workorder/<wo\_id>/solutions & Append a solution-log entry. \\
GET / POST & /api/workorder/<wo\_id>/attachments & List / upload attachments. \\
GET & /api/attachments/<att\_id> & Download an attachment. \\
GET & /api/workorder/<wo\_id>/audit & WO change history (admin-gated). \\
POST & /api/workorder/<wo\_id>/assign & Assign to a technician (Sean-gated). \\
\bottomrule
\end{longtable}

\subsubsection{Structure Management}
\begin{longtable}{@{}p{0.17\textwidth}>{\ttfamily\small}p{0.44\textwidth}p{0.31\textwidth}@{}}
\toprule
\textbf{Method} & \textbf{Path} & \textbf{Description} \\
\midrule
\endhead
POST & /api/departments & Create a department. \\
POST & /api/machines & Create a machine. \\
POST & /api/departments/<dk>/deactivate\ \textbar\ restore & Soft-delete / restore a department. \\
POST & /api/machines/deactivate\ \textbar\ restore & Soft-delete / restore a machine. \\
POST & /api/departments/<dk>/machines/<eq>/nickname & Set/clear a display nickname (admin-gated). \\
GET / POST & /api/technicians & List / add technicians. \\
DELETE & /api/technicians/<name> & Remove a technician. \\
GET & /api/roles, /api/roles/views & List active roles; public role/view-key pairs. \\
POST & /api/roles & Add a role with aliases (Sean-gated). \\
DELETE & /api/roles/<name> & Delete a role (Sean-gated). \\
POST & /api/roles/<name>/view & Assign a view profile to a role (master-gated). \\
GET / POST / DELETE & /api/view-profiles[/<key>] & List (open) / create-update / delete (master-gated) view profiles. \\
GET / POST & /api/departments/<dk>/floorplan & Read / save (admin-gated) a department floor plan. \\
\bottomrule
\end{longtable}

\subsubsection{Spare Parts, Machine Info, and Vendor Contacts}
\begin{longtable}{@{}p{0.19\textwidth}>{\ttfamily\small}p{0.43\textwidth}p{0.30\textwidth}@{}}
\toprule
\textbf{Method} & \textbf{Path} & \textbf{Description} \\
\midrule
\endhead
GET & /api/spare-parts & Global spare-parts list, newest first. \\
GET & /api/spare-parts/options & Departments/machines/divisions for the create/edit form. \\
POST / PATCH / DELETE & /api/spare-parts[/<id>] & Create / edit / delete a part (admin-gated). \\
GET & /api/departments/<dk>/machines/<eq>/spare-parts & Parts linked to one machine. \\
GET, PATCH / POST & /api/departments/<dk>/machines/<eq>/info & Read / update the editable Machine Info record. \\
POST & /api/departments/<dk>/machines/<eq>/parse-notes & LLM sorts free-form notes into structured contacts/logins/purchase info. \\
GET & /api/vendor-contacts & Combined vendor list. \\
POST / PATCH / DELETE & /api/vendor-contacts[/<id>] & Create / edit / delete vendor contacts (edit-gated). \\
POST / DELETE & /api/vendor-contact-types[/<name>] & Manage contact types (edit-gated). \\
\bottomrule
\end{longtable}

\subsubsection{Analytics and AI}
\begin{longtable}{@{}p{0.17\textwidth}>{\ttfamily\small}p{0.46\textwidth}p{0.29\textwidth}@{}}
\toprule
\textbf{Method} & \textbf{Path} & \textbf{Description} \\
\midrule
\endhead
GET & /api/departments/<dk>/machines/<eq>/trends & Per-machine monthly KPI trends. \\
GET & /api/departments/<dk>/trends & Department monthly trends. \\
GET & /api/departments/<dk>/machines-trends & Per-machine yearly totals (Pareto chart). \\
GET & \ldots/month-synopsis & LLM synopsis of what drove a month's stats. \\
GET / POST / PUT & \ldots/guide & Read / generate / save-edit the checklist (edit-gated for writes). \\
POST & \ldots/guide/update & AI update with new work orders. \\
GET & \ldots/guide/versions[/<n>] & Version history / a specific version. \\
POST & \ldots/guide/versions/<n>/restore & Restore an old checklist version. \\
GET / POST / DELETE & /api/calendar-events[/<id>] & Calendar entries. \\
GET / POST / DELETE & /api/chart-events[/<id>] & Global chart timeline markers (edit-gated). \\
\bottomrule
\end{longtable}

\subsubsection{System and E-mail}
\begin{longtable}{@{}p{0.17\textwidth}>{\ttfamily\small}p{0.44\textwidth}p{0.31\textwidth}@{}}
\toprule
\textbf{Method} & \textbf{Path} & \textbf{Description} \\
\midrule
\endhead
POST & /api/reload & Force a data reload. \\
GET & /api/edit-auth, /api/admin-auth, /api/sean-auth & Whether each gate is password-protected. \\
POST & /api/verify-\{edit\textbar admin\textbar sean\}-password & Verify a supplied password. \\
GET / POST & /api/email/status, /api/email/recipients & E-mail configuration status / recipient list. \\
POST & /api/email/test, /api/email/weekly-digest & Send a test e-mail / trigger the weekly digest. \\
\bottomrule
\end{longtable}

% =============================================================================
\section{Intelligence Layer (LLM Integration)}
% =============================================================================

\subsection{Providers and Transport}
Two providers are supported, selected by environment:
\begin{itemize}
    \item \textbf{Ollama Cloud} (default when \code{OLLAMA\_API\_KEY} is
          set).
    \item \textbf{Google Gemini} (\code{GEMINI\_API\_KEY}).
\end{itemize}
Because the corporate network blocks the standard Python TLS stack, HTTPS
calls are made through a Windows-native PowerShell transport instead. The
module also computes local statistics (material cost, labor time, downtime
hours) and builds prompts.

\subsection{Troubleshooting-Checklist Engine}
Single source of truth for creating, editing, archiving, and regenerating
per-machine checklists, used both interactively and by the nightly job.

\textbf{Manual-edit preservation.} When an operator saves a hand-edited
checklist, the previous file is backed up, the old/new Markdown is diffed,
and the added/removed lines are appended to an edit log. That log is
rendered into an ``operator edits that must be preserved'' block injected
into every future prompt, so edits survive regeneration. Entries are
applied oldest-first with later wins, so edits-of-edits are handled
automatically.

% =============================================================================
\section{E-mail Notifications}
% =============================================================================

Three flows, each rendered as HTML with a plain-text fallback:
\begin{enumerate}
    \item \textbf{New unscheduled WO} --- sent immediately on creation.
    \item \textbf{Closed and completed} --- sent when a WO transitions into
          a completed status (transition-only; re-saving does not
          re-notify).
    \item \textbf{Weekly digest} --- scheduled work due this coming week
          plus any past-due scheduled work still pending, sent by a daemon
          thread on a configurable schedule.
\end{enumerate}

Transport is plain SMTP, so a self-serve Gmail App Password, Office 365, or
any SMTP server works. The module is intentionally inert until credentials
are configured and at least one recipient is set.

% =============================================================================
\section{Presentation Layer (Frontend)}
% =============================================================================

The entire UI is a single self-contained file, \code{webapp/index.html},
served by Flask with no-cache headers. Technology choices:
\begin{itemize}
    \item \textbf{React} with \textbf{htm} tagged templates and
          \textbf{Babel standalone} for in-browser compilation --- no build
          step.
    \item \textbf{Tailwind CSS} for styling.
    \item \textbf{marked} to render checklist Markdown; \textbf{jsPDF} for
          client-side PDF export.
    \item All libraries vendored locally: works fully offline on the LAN.
\end{itemize}

Client-side routing uses a view state of company $\rightarrow$ division
$\rightarrow$ department $\rightarrow$ machine, plus an inactive view
listing soft-deleted items with restore actions. The machine page has tabs
for Dashboard, AI Troubleshooting Checklist, Unscheduled WOs, and Scheduled
WOs, plus Machine Info and trend/Pareto charts. The Unscheduled/Scheduled WO
tabs include CSV/JSON and client-side PDF export buttons for a machine's
work-order list.

\textbf{Nicknames.} Any machine can be given an admin-set display nickname
while the real PM equipment name is preserved and still shown; grouping and
statistics continue to key off the real PM name.

\textbf{Floor Plan.} Selected departments have an interactive,
percent-positioned Floor Plan tab: a sidebar lists linked machines with live
stats and a highlight for any with a critical open WO; admin edit mode
supports dragging, resizing, linking/relabeling, and adding boxes.

\textbf{Spare Parts.} A division-scoped spare-parts inventory screen lists
parts with description, type, location, quantity, condition, linked
machine(s), and purchasing info; a part can also be viewed from the machine
it is linked to.

\textbf{Roles and View Profiles.} Beyond individually tracked technicians,
shared roles (e.g.\ ``QA Technician'') let multiple people sign in as a team
identity for My Schedule. A master-gated admin page defines named view
profiles (curated sets of enabled features) and assigns one to each role,
so a role's users only see the subset of MINT relevant to them.

% =============================================================================
\section{Configuration Reference}
\label{sec:config}
% =============================================================================

All configuration is via environment variables, typically in a git-ignored
\code{.env} file.

\begin{longtable}{@{}>{\ttfamily}p{0.30\textwidth}>{\ttfamily}p{0.14\textwidth}p{0.48\textwidth}@{}}
\toprule
\textbf{Variable} & \textbf{Default} & \textbf{Meaning} \\
\midrule
\endhead
SERVER\_HOST / SERVER\_PORT & 0.0.0.0 / 5000 & Bind address and port. \\
EDIT\_PASSWORD & (blank) & Checklist/vendor/chart-event edit gate; blank = open. \\
ADMIN\_PASSWORD & (blank) & Delete-WO / history / spare-parts / nickname / floor-plan gate; blank = open. \\
SEAN\_PASSWORD & sean & Assignment + team-dashboard gate. \\
MASTER\_PASSWORD & mintmaster & View-profile / role-view management gate. \\
OLLAMA\_API\_KEY / OLLAMA\_MODEL & -- / gpt-oss:120b & Ollama Cloud checklist provider. \\
GEMINI\_API\_KEY / GEMINI\_MODEL & -- / gemini-2.5-flash & Gemini provider (fallback). \\
CHROME\_AUTOSTART & 1 & 0 = do not auto-launch debug Chrome. \\
CHROME\_DEBUG\_PORT & 0 & 0 = auto-pick and capture; or pin a port. \\
CHROME\_PATH / CHROME\_PROFILE & auto & Non-standard Chrome binary / profile path. \\
NIGHTLY\_DEPT\_TIMEOUT & 14400 & Per-department scrape timeout (seconds). \\
EMAIL\_ENABLED & auto & 1/0; defaults on when SMTP credentials exist. \\
SMTP\_HOST / SMTP\_PORT & -- / 587 & Mail server. \\
SMTP\_USER / SMTP\_PASSWORD / SMTP\_FROM & -- & Credentials; FROM defaults to USER. \\
SMTP\_STARTTLS & 1 & 0 disables STARTTLS. \\
GMAIL\_ADDRESS / GMAIL\_APP\_PASSWORD & -- & Gmail shortcuts for the SMTP\_* vars. \\
MINT\_BASE\_URL & http://localhost:5000 & Base URL used in e-mail links. \\
WEEKLY\_DIGEST\_ENABLED / DAY / HOUR & 1 / 0 / 7 & Weekly digest schedule (0 = Monday). \\
\bottomrule
\end{longtable}

% =============================================================================
\section{Setup and Deployment}
% =============================================================================

\subsection{Prerequisites}
\begin{itemize}
    \item Windows machine (the LLM transport and Chrome management are
          Windows-specific), Python 3.10+, Google Chrome installed.
    \item Python dependencies from \code{requirements.txt}: selenium,
          webdriver-manager, beautifulsoup4, lxml, google-genai,
          python-dotenv, flask, flask-cors, APScheduler.
\end{itemize}

\subsection{One-Time Setup}
\begin{lstlisting}
pip install -r requirements.txt

# Create .env with real keys:
#   GEMINI_API_KEY=...
#   OLLAMA_API_KEY=...
#   EDIT_PASSWORD=...  ADMIN_PASSWORD=...  (optional)
\end{lstlisting}
\code{.env} is git-ignored and must never be committed.

\subsection{Chrome Login Session}
\begin{lstlisting}
python server.py          # auto-launches debug Chrome, captures the port
# First time: log into the PM site in that Chrome window; leave it open.

python chrome_session.py            # manual launch + port
python chrome_session.py --status   # running port + login state
\end{lstlisting}
If the session expires, scrapes are skipped (never overwriting good data)
and the failure is reported.

\subsection{Running the Server}
\begin{lstlisting}
python server.py
#   Local:   http://127.0.0.1:5000
#   Network: http://<host-ip>:5000   (LAN)
\end{lstlisting}
For local development/testing without Chrome or the nightly machinery:
\begin{lstlisting}
$env:CHROME_AUTOSTART="0"; $env:NIGHTLY_ENABLED="0"; python server.py
\end{lstlisting}

\subsection{LAN Access}
The server binds \code{0.0.0.0}; open the firewall once (Admin PowerShell):
\begin{lstlisting}
New-NetFirewallRule -DisplayName "PM Dashboard" -Direction Inbound `
    -Action Allow -Protocol TCP -LocalPort 5000
\end{lstlisting}

\subsection{Manual Data Operations}
\begin{lstlisting}
python scraper.py --limit 3      # quick scrape smoke test
python scraper.py --from-html    # offline replay from ./pages/
python run_parallel.py --division bed --refresh-profiles
python nightly_update.py         # full rescrape + regeneration cycle
python analyze_equipment.py --equipment-id 1877   # ad-hoc LLM analysis
curl -X POST http://localhost:5000/api/reload      # hot-reload data
\end{lstlisting}

\subsection{Operational Notes}
\begin{itemize}
    \item \textbf{Stale server process:} a previously started
          \code{server.py} can hold port 5000 so a new instance silently
          fails to bind while the old code keeps answering. Stop the old
          \code{python.exe} running \code{server.py} before restarting.
    \item \textbf{Attachments and databases} live under \code{mint\_data/}
          and \code{guides/} --- include both in any backup strategy.
    \item \textbf{Division scraping} requires re-selecting the target
          division in the debug Chrome dropdown before each division's
          run, otherwise the session silently keeps scraping the previous
          division.
\end{itemize}

% =============================================================================
\section{Security Considerations}
% =============================================================================

\begin{itemize}
    \item The application is designed for a trusted LAN: reads are
          unauthenticated; writes are gated by shared passwords transmitted
          in headers over plain HTTP. It must not be exposed to the
          internet without a reverse proxy adding TLS and real
          authentication.
    \item Scraping reuses an interactive Chrome login --- no PM credentials
          are stored anywhere in the codebase.
    \item Secrets live only in the git-ignored \code{.env}. Any API keys
          that were historically committed must be treated as leaked and
          rotated.
    \item Every mutation is attributable via the mandatory author name and
          the audit log; scraped records are immutable (edits are overlay
          patches) and manual WOs are the only deletable records.
\end{itemize}

% =============================================================================
\section{Known Limitations and Future Work}
% =============================================================================

\begin{itemize}
    \item Single-process, single-host deployment; SQLite and in-memory
          caches limit horizontal scaling (acceptable for a facility-level
          tool).
    \item The Windows PowerShell TLS transport ties LLM calls to Windows.
    \item Author identity is honor-system (a typed name), not authenticated
          accounts; the same is true for shared role logins.
    \item BLA and BED are fully populated; other divisions (BMC, BWEC, KKP,
          BGD, CIT) are scraped incrementally with the same division-aware
          pipeline described in Section~\ref{sec:divisions}.
    \item The frontend is one large HTML file; a build pipeline with proper
          module splitting would improve maintainability.
    \item The e-mail digest scheduler uses a polling daemon thread;
          migrating to APScheduler jobs (already a dependency) would be
          cleaner.
\end{itemize}

\end{document}
